\subsection{2.12 Aktueller Stand der numerischen Evidenz und offene Forschungsfragen}

Nach Abschluss der bisherigen numerischen Untersuchungen lässt sich erstmals eine konsistente Zwischenbilanz ziehen.

Die zahlreichen Simulationen verfolgten ursprünglich kein einheitliches Ziel. Erst im Verlauf der Arbeiten kristallisierte sich Schritt für Schritt heraus, welche Fragestellungen tatsächlich von fundamentaler Bedeutung sind.

Im Rückblick lassen sich die Ergebnisse in drei Kategorien einteilen:

\begin{enumerate}
\item robuste numerische Befunde,
\item wahrscheinliche Hypothesen,
\item weiterhin offene Probleme.
\end{enumerate}

%%%%%%%%%%%%%%%%%%%%%%%%%%%%%%%%%%%%%%%%%%%%%%%%%%%%%%%%%%%%%%

\section*{I. Robuste numerische Befunde}

Mehrere Resultate wurden unabhängig voneinander reproduziert und besitzen daher einen deutlich höheren Evidenzgrad.

%%%%%%%%%%%%%%%%%%%%%%%%%%%%%%%%%%%%%%%%%%%%%%%%%%%%%%%%%%%%%%

\subsubsection*{(1) Der rohe OPP ist hochdimensional}

Ohne jede Projektion wächst die effektive Dimension nahezu linear mit der Anzahl der Attraktoren.

Typische Messwerte waren

\[
\begin{array}{c|c}
N & d_{\mathrm{eff}} \\
\hline
10 & 9.6 \\
20 & 17.8 \\
40 & 31.6
\end{array}
\]

Daraus folgt unmittelbar:

\[
\boxed{
\text{Der OPP besitzt keine intrinsische vierdimensionale Struktur.}
}
\]

Die Vier muss daher zwangsläufig erst durch einen weiteren Mechanismus entstehen.

%%%%%%%%%%%%%%%%%%%%%%%%%%%%%%%%%%%%%%%%%%%%%%%%%%%%%%%%%%%%%%

\subsubsection*{(2) Projektionen verändern die effektive Dimension massiv}

Alle getesteten Projektionsoperatoren reduzierten die effektive Dimension.

Allerdings nicht in gleichem Maße.

Die Projektion ist daher kein nebensächlicher mathematischer Schritt,

sondern bestimmt unmittelbar die beobachtbare Geometrie.

%%%%%%%%%%%%%%%%%%%%%%%%%%%%%%%%%%%%%%%%%%%%%%%%%%%%%%%%%%%%%%

\subsubsection*{(3) Nicht jede Projektion ist physikalisch sinnvoll}

Der Vergleich verschiedener Verfahren ergab deutliche Qualitätsunterschiede.

Die Rangfolge lautete ungefähr:

\[
\begin{array}{l|c}
\text{Methode}
&
d_{\mathrm{eff}}
\\
\hline
\text{Random Projection}
&
2.9
\\
\text{PCA}
&
3.6
\\
\text{Kernel PCA}
&
3.9
\\
\text{Diffusion Maps}
&
4.0
\end{array}
\]

Bereits dieses Resultat spricht gegen die Annahme,

jede beliebige Projektion könne eine physikalische Raumzeit erzeugen.

%%%%%%%%%%%%%%%%%%%%%%%%%%%%%%%%%%%%%%%%%%%%%%%%%%%%%%%%%%%%%%

\subsubsection*{(4) Diffusionsprojektionen besitzen einen stabilen Arbeitsbereich}

Die umfangreichen Parameterstudien zeigten,

dass sich

\[
d_{\mathrm{eff}}
\approx
4
\]

nur innerhalb eines begrenzten Bereiches der Diffusionszeit ergibt.

Insbesondere

\[
t
\approx
9
-
10
\]

lieferte wiederholt Werte nahe vier.

%%%%%%%%%%%%%%%%%%%%%%%%%%%%%%%%%%%%%%%%%%%%%%%%%%%%%%%%%%%%%%

\section*{II. Wahrscheinliche Hypothesen}

Mehrere Schlussfolgerungen werden durch die bisherigen Daten gestützt,

sind jedoch noch nicht endgültig bewiesen.

%%%%%%%%%%%%%%%%%%%%%%%%%%%%%%%%%%%%%%%%%%%%%%%%%%%%%%%%%%%%%%

\subsubsection*{(1) Raumzeit ist eine Projektion}

Die numerischen Ergebnisse sprechen deutlich dafür,

dass Raumzeit keine Eigenschaft des OPP selbst ist.

Vielmehr scheint zu gelten

\[
\boxed{
\text{Raumzeit}
=
\mathcal{P}(\mathrm{OPP}).
}
\]

Diese Aussage bildet inzwischen eine der zentralen Arbeitshypothesen der PST.

%%%%%%%%%%%%%%%%%%%%%%%%%%%%%%%%%%%%%%%%%%%%%%%%%%%%%%%%%%%%%%

\subsubsection*{(2) Projektionen entsprechen Beobachtungsprozessen}

Die Ähnlichkeit zwischen Informationsreduktion

und realer Beobachtung wurde im Verlauf der Arbeiten immer deutlicher.

Dadurch entstand die Vermutung,

dass physikalische Beobachter mathematisch selbst als Projektionsoperatoren beschrieben werden können.

%%%%%%%%%%%%%%%%%%%%%%%%%%%%%%%%%%%%%%%%%%%%%%%%%%%%%%%%%%%%%%

\subsubsection*{(3) Symmetriebruch könnte projektionsbedingt sein}

Die Diskussion über Eichgruppen,

Lie-Gruppen

und Kaluza-Klein-Strukturen führte zu einer neuen Interpretation.

Nicht der OPP muss bereits alle beobachteten Symmetrien besitzen.

Vielmehr könnte gelten:

\[
\boxed{
\text{Symmetriebruch}
=
\text{Folge der Projektion.}
}
\]

Dies stellt derzeit eine der interessantesten offenen Forschungsrichtungen dar.

%%%%%%%%%%%%%%%%%%%%%%%%%%%%%%%%%%%%%%%%%%%%%%%%%%%%%%%%%%%%%%

\subsubsection*{(4) Fake-5D}

Die häufig beobachteten Werte

\[
d_{\mathrm{eff}}
\approx
5
\]

werden inzwischen nicht mehr als Hinweis auf eine reale fünfte Raumdimension interpretiert.

Wahrscheinlicher erscheint,

dass eine zusätzliche Spektralmode bestehen bleibt,

welche beispielsweise

\begin{itemize}

\item

Projektionsunschärfe,

\item

innere Freiheitsgrade,

\item

oder Quantenfluktuationen

repräsentiert.

\end{itemize}

%%%%%%%%%%%%%%%%%%%%%%%%%%%%%%%%%%%%%%%%%%%%%%%%%%%%%%%%%%%%%%

\section*{III. Offene Probleme}

Trotz der bisherigen Fortschritte bleiben zahlreiche fundamentale Fragen unbeantwortet.

%%%%%%%%%%%%%%%%%%%%%%%%%%%%%%%%%%%%%%%%%%%%%%%%%%%%%%%%%%%%%%

\subsubsection*{(1) Der Projektionsoperator}

Bislang existiert kein physikalisch motivierter Projektionsoperator.

Kernel PCA und Diffusion Maps sind lediglich mathematische Kandidaten.

Die eigentliche Projektion der Natur ist weiterhin unbekannt.

%%%%%%%%%%%%%%%%%%%%%%%%%%%%%%%%%%%%%%%%%%%%%%%%%%%%%%%%%%%%%%

\subsubsection*{(2) Lorentz-Metrik}

Bis heute wurde nicht gezeigt,

wie aus einer Projektion zwingend eine Lorentz-Signatur

\[
(-,+,+,+)
\]

entsteht.

Dieses Problem besitzt höchste Priorität.

%%%%%%%%%%%%%%%%%%%%%%%%%%%%%%%%%%%%%%%%%%%%%%%%%%%%%%%%%%%%%%

\subsubsection*{(3) Eichgruppen}

Ebenso offen bleibt,

ob sich

\[
SU(3)
\times
SU(2)
\times
U(1)
\]

tatsächlich aus projektionsbedingten Restgruppen ableiten lassen.

%%%%%%%%%%%%%%%%%%%%%%%%%%%%%%%%%%%%%%%%%%%%%%%%%%%%%%%%%%%%%%

\subsubsection*{(4) Dynamik}

Die bisherigen Untersuchungen behandeln überwiegend statische Strukturen.

Eine vollständige Theorie benötigt jedoch zusätzlich

\begin{itemize}

\item

Zeitentwicklung,

\item

Feldgleichungen,

\item

und Wechselwirkungen.

\end{itemize}

%%%%%%%%%%%%%%%%%%%%%%%%%%%%%%%%%%%%%%%%%%%%%%%%%%%%%%%%%%%%%%

\subsubsection*{(5) Mathematische Fundierung}

Viele bisher verwendete Größen

wie beispielsweise

\[
d_{\mathrm{eff}},
\]

das Stabilitätsfunktional,

oder verschiedene Projektionsoperatoren,

wurden zunächst heuristisch eingeführt.

Eine streng mathematische Herleitung steht noch aus.

%%%%%%%%%%%%%%%%%%%%%%%%%%%%%%%%%%%%%%%%%%%%%%%%%%%%%%%%%%%%%%

\section*{Der wichtigste Erkenntnisgewinn}

Trotz aller offenen Fragen besitzt die bisherige Arbeit einen entscheidenden Erfolg.

Zu Beginn der Untersuchungen bestand die Hoffnung,

im OPP direkt eine vierdimensionale Struktur zu finden.

Nach zahlreichen Simulationen lautet die heutige Sichtweise genau umgekehrt.

\[
\boxed{
\text{Vierdimensionale Raumzeit muss gar nicht im OPP existieren.}
}
\]

Sie entsteht vielmehr erst durch geeignete Projektionsprozesse.

Dieser Perspektivwechsel verändert die gesamte Fragestellung der PST.

%%%%%%%%%%%%%%%%%%%%%%%%%%%%%%%%%%%%%%%%%%%%%%%%%%%%%%%%%%%%%%

\section*{Neue Forschungsstrategie}

Die zukünftige Entwicklung der PST orientiert sich daher an folgendem Schema:

\[
\boxed{
\text{Existenz}
\rightarrow
\text{OPP}
\rightarrow
\text{Projektionsoperator}
\rightarrow
\text{Symmetriebruch}
\rightarrow
\text{Raumzeit}
\rightarrow
\text{Physik}.
}
\]

Nicht mehr die Geometrie des OPP steht im Mittelpunkt,

sondern die Mathematik der Projektion.

%%%%%%%%%%%%%%%%%%%%%%%%%%%%%%%%%%%%%%%%%%%%%%%%%%%%%%%%%%%%%%

\section*{Schlussbemerkung}

Die numerischen Arbeiten der ersten Entwicklungsphase haben weder die PST bewiesen noch widerlegt.

Sie haben jedoch etwas erreicht,

das wissenschaftlich mindestens ebenso wichtig ist.

Sie haben die ursprüngliche Fragestellung präzisiert.

Heute ist wesentlich klarer,

\begin{itemize}

\item

welche Hypothesen mit den Daten vereinbar sind,

\item

welche Ideen aufgegeben wurden,

\item

welche Projektionsverfahren Erfolg versprechen,

\item

und an welchen mathematischen Problemen künftig gearbeitet werden muss.

\end{itemize}

Damit endet der erste große numerische Entwicklungszyklus der PST.

Die weitere Forschung richtet ihren Schwerpunkt nun auf die mathematische Rekonstruktion derjenigen Projektion,

aus der die bekannte Physik als stabile emergente Struktur hervorgeht.
